\begin{solapartat}
\punts{0,3} Segons la tercera llei de Kepler:
\[ \frac{a^3}{T^2} = \frac{G M_\mathrm{Sol}}{(2\pi)^2} \]
On $a$ és la longitud del semieix major.

\punts{0,3}
\begin{align*}
a &= \sqrt[3]{G M_\mathrm{Sol}\left(\frac{T}{2\pi}\right)^2} = \sqrt[3]{6,67\times 10^{-11}\times 1,99\times 10^{30}\left(\frac{76\times 365\times 24\times 3600}{2\pi}\right)^2}\\
&= 2,68\times 10^{12}\ \mathrm{m}
\end{align*}

\destaca{Alternativament}, s'ha d'acceptar com correcte que l'estudiant recordi que la tercera llei de Kepler al sistema solar s'expressa com
\[ \frac{a^3}{T^2} = 1\ \frac{\mathrm{ua}^3}{\mathrm{any}^2} \]
Llavors,
\[ a = \sqrt[3]{T^2} = \sqrt[3]{76^2} = 17,94\ \mathrm{ua} \]

\punts{0,65}
\figura[0.5]{sol_fig1.png}

Cal dibuixar una òrbita clarament el·líptica amb dos focus.

Cal dibuixar el sol en un dels focus, sinó es resta 0,3 p.

Cal dibuixar el periheli i l'afeli sobre l'eix major, un a cada extrem amb el periheli sent el més proper al sol, sinó es resta 0,3 p.
\end{solapartat}

\begin{solapartat}
\punts{0,1} Les distàncies a l'afeli i al periheli sumades són igual a la longitud del eix major:
\[ r_A + r_P = 2a \]

$r_P = 0,586\ \mathrm{ua}\,\dfrac{1,50\times 10^{11}\ \mathrm{m}}{1\ \mathrm{UA}} = 8,79\times 10^{10}\ \mathrm{m}$ \punts{0,1}

$r_A = 2a - r_P = 5,28\times 10^{12}\ \mathrm{m}$ \punts{0,1}

\punts{0,3} Segons la llei de gravitació universal, el mòdul de la força s'expressa com:
\[ F = G\,\frac{m_\mathrm{Halley} M_\mathrm{Sol}}{r^2} \]
i $F = m_\mathrm{Halley}g$, llavors:

\punts{0,25} $g = G\,\dfrac{M_\mathrm{Sol}}{r^2}$

Per l'afeli: $g = G\,\dfrac{M_\mathrm{Sol}}{r_A^2} = 4,76\times 10^{-6}\ \mathrm{m/s^2}$ \punts{0,2}

I pel periheli: $g = G\,\dfrac{M_\mathrm{Sol}}{r_A^2} = 1,72\times 10^{-2}\ \mathrm{m/s^2}$ \punts{0,2}
% DUBTE: a l'original, la fórmula del camp gravitatori al periheli porta per error el subíndex
% "A" (d'afeli) al denominador en lloc de "P" (de periheli); el resultat numèric, 1,72e-2 m/s^2,
% sí que correspon a r_P. Es transcriu tal com surt a la pauta oficial.
\end{solapartat}
